\chapter{Numerical Optimization}
\label{ch:optimization}

Section~\ref{ssec:optimization} gave a short tour of optimization in Mathilda.
It minimized Rosenbrock's function, solved one constrained problem, and let
differential evolution find the bottom of Rastrigin's egg carton. This chapter
goes further. It is about the algorithms behind the four optimizers
\B{FindMinimum}, \B{FindMaximum}, \B{NMinimize}, and \B{NMaximize}: what each
one computes, why it works, when it fails, and how to tell the difference at
the prompt.

Optimization asks for the smallest (or largest) value of a function, and for the
point where it occurs. That single question covers a surprising amount of
applied mathematics. Fitting a model to data is the minimization of a sum of
squares. The shape of a hanging chain minimizes potential energy. A light ray
takes the path of least time. A factory schedule maximizes profit subject to
limited machine hours. A portfolio minimizes risk for a required return. The
last section of the chapter solves each of these as a short worked example,
with a picture of the answer.

Three themes run through the chapter. The first is the mathematics. Every
method is introduced with the idea that makes it work. Brent's method rests on
the golden ratio, BFGS on the secant condition, differential evolution on its
mutation rule, and the penalty method on hidden Lagrange multipliers. The second is
honesty about what a numerical optimizer can promise. A local method finds a
local minimum. A global method searches harder but proves nothing. A
constrained answer must satisfy its constraints, and Mathilda refuses to return
one that does not. The third is speed. The optimizers compile the objective to
bytecode before they search, and the benchmarks at the end of the chapter show
what that buys against SciPy and Mathematica, and where it does not help.

\section{Local and global}
\label{sec:opt-local-global}

A \emph{local minimum}\index{local minimum} of \(f\) is a point \(x^*\) with
\(f(x^*) \le f(x)\) for every \(x\) in some neighbourhood of \(x^*\). A
\emph{global minimum}\index{global minimum} satisfies the inequality for every
\(x\) in the whole domain. A function can have many local minima and still
have only one global minimum. The function \(x\cos x\) on \([0, 16]\) is a
good first example. It oscillates with growing amplitude, so each valley is
deeper than the one before.

\plotcell{optimization/landscape}

The red points are the three answers \B{FindMinimum} gives from the starting
points \(2\), \(7\), and \(14\). Each one is correct. Each is the bottom of the
valley that contains its starting point. None of them knows about the others.
This is the defining property of a \emph{local} method. It follows the slope
downhill from where you put it, and it stops at the first point where the slope
vanishes.

The global optimizers take a region instead of a starting point.
\B{NMinimize} searches the whole interval and reports the deepest valley it
finds:

\begin{codepairs}
\pair{optimization/forms/1}{From \(x = 2\), \B{FindMinimum} finds the first
valley.}
\pair{optimization/forms/2}{From \(x = 7\) it goes downhill to the right and
finds the second.}
\pair{optimization/forms/3}{\B{NMinimize} searches all of \([0,16]\). Its
answer is the deepest of the three valleys in the figure.}
\pair{optimization/forms/4}{\B{FindMaximum} climbs uphill. From \(x=5\) it
reaches the nearest peak.}
\pair{optimization/forms/5}{\B{NMaximize} finds the highest peak in the
interval.}
\pair{optimization/forms/6}{Every optimizer returns a list of two things. The
first is the optimal value and the second is a list of rules for the optimal
point.}
\pair{optimization/forms/7}{The rules plug straight into \mcode{ReplaceAll}.}
\pair{optimization/forms/8}{So it is easy to check the answer by
substituting it back into the objective.}
\end{codepairs}

The two families use the same syntax as far as possible. The first argument is
the objective, or a list \mcode{\{f, cons\}} of an objective and its
constraints. The second argument names the variables. \B{FindMinimum} pairs
each variable with a starting value, as in \mcode{\{x, 2\}}. \B{NMinimize}
takes the bare variables and searches a region. Both return
\mcode{\{fmin, \{x -> xmin, \ldots\}\}}.

\calloutnote{Theory}{There is no general algorithm that finds the global
minimum of an arbitrary continuous function using finitely many function
values. Between any two samples the function could dip arbitrarily low, and no
finite set of samples can rule that out. A global method therefore needs extra
knowledge, such as a Lipschitz bound, or it must accept a probabilistic
answer. The methods in \B{NMinimize} take the second route, except DIRECT,
which uses the Lipschitz idea without knowing the constant.}

\begin{table}[ht]
\centering
\small
\begin{tabular}{@{}lll@{}}
\toprule
& local (needs a start) & global (needs a region) \\
\midrule
minimize & \mcode{FindMinimum} & \mcode{NMinimize} \\
maximize & \mcode{FindMaximum} & \mcode{NMaximize} \\
default method & Brent (1 variable), BFGS (several) & differential evolution \\
constraints & boxes, smooth inequalities and equalities & also \mcode{Or} and integers \\
answer & a local optimum & the best optimum found \\
\bottomrule
\end{tabular}
\caption{The four optimizers.}
\label{tab:opt-four}
\end{table}

The maximizers are thin wrappers. \B{FindMaximum}\mcode{[f, ...]} calls
\B{FindMinimum}\mcode{[-f, ...]} and negates the value it gets back, and
\B{NMaximize} does the same with \B{NMinimize}. Everything this chapter says
about minimization therefore applies to maximization with the sign reversed.

\section{Minimizing in one variable}
\label{sec:opt-brent}

With one variable and no derivative, the natural idea is to shrink an interval
that is known to contain a minimum. Suppose \(a < b < c\) and \(f(b)\) is
smaller than both \(f(a)\) and \(f(c)\). A continuous \(f\) must have a local
minimum in \((a, c)\). Such a triple is a \emph{bracket}\index{bracket (minimization)}.
Evaluate \(f\) at a new point \(u\) inside the larger of the two subintervals.
Whichever of \(b\) and \(u\) has the smaller value becomes the middle of a new,
smaller bracket.

The best place for \(u\) is not the midpoint. If the new point divides the
larger subinterval in the ratio
\[
  \frac{3 - \sqrt5}{2} \approx 0.382,
\]
then the new bracket has the same proportions as the old one, and every step
shrinks the interval by the same factor \(1/\varphi \approx 0.618\), where
\(\varphi\) is the golden ratio. This is \emph{golden-section
search}\index{golden-section search}. It is slow but it never fails, because it
needs only comparisons of function values.

Near a smooth minimum a function looks like a parabola, and a parabola through
three points can be minimized exactly. \emph{Brent's
method}\index{Brent's method} combines the two ideas. It fits a parabola through
the three best points and jumps to its vertex. When that jump is unsafe, because
it lands outside the bracket or the steps stop shrinking fast enough, Brent's
method takes a golden-section step instead. The result converges superlinearly
on smooth functions and never does worse than golden section on rough ones.
It is the default method for \B{FindMinimum} in one variable.

\begin{codepairs}
\pair{optimization/brent/1}{The quartic \(x^4 - 3x^2 + x\) has two minima.
From \(x = 1\), Brent finds the shallow one.}
\pair{optimization/brent/2}{From \(x=-1\) it finds the deeper one.}
\pair{optimization/brent/3}{The form \mcode{\{x, x0, xmin, xmax\}} gives
the search interval explicitly.}
\pair{optimization/brent/4}{Brent needs no derivative, so a kink is no obstacle.
The minimum of \(|x - 1/3|\) is found to the default tolerance.}
\pair{optimization/brent/5}{Twelve evaluations of the objective are enough for
a smooth minimum. \mcode{EvaluationMonitor} counts them.}
\end{codepairs}

When \B{FindMinimum} is given a single starting value, it first builds a
bracket by stepping downhill with growing steps until the function turns up
again. That search is the reason a starting point matters. The bracket is built
around the first valley the steps reach.

\calloutnote{Under the hood}{The one-dimensional path is \mcode{fm\_bracket}
followed by \mcode{fm\_brent\_min} in
\texttt{src/numerical\_calculus/findmin\_common.c}. Both have MPFR twins
(\mcode{fm\_bracket\_mpfr} and \mcode{fm\_brent\_min\_mpfr}) with the same
acceptance test and the same golden-section fallback, so
\mcode{WorkingPrecision} runs the identical algorithm in higher precision.}

\subsection{How many digits can a minimum have?}

The minimum of the Gamma function on the positive axis is a classical
constant. \B{FindMinimum} finds it at machine precision and, with
\mcode{WorkingPrecision}, in 40-digit arithmetic. The minimum is also a root of
the derivative \(\psi(x) = \Gamma'(x)/\Gamma(x)\), so \B{FindRoot} applied to
\B{PolyGamma} gives an independent check:

\begin{codepairs}
\pair{optimization/precision/1}{The minimum of \(\Gamma\) at machine
precision.}
\pair{optimization/precision/2}{The same minimum in 40-digit arithmetic. The
value is printed to 40 digits.}
\pair{optimization/precision/3}{The root of \(\psi\), computed by
\B{FindRoot}.}
\end{codepairs}

Compare the two locations digit by digit. They agree for the first 21
significant digits and then part company, even though both computations
carried 40 digits. The minimizer is not wrong. It has reached the limit that
every minimizer faces.

\begin{theorem}[Accuracy of a minimizer]
\index{minimizer!attainable accuracy}
Let \(f\) be twice differentiable with \(f'(x^*) = 0\) and \(f''(x^*) > 0\).
If \(f\) can be evaluated only to relative precision \(\varepsilon\), then
\(x^*\) can be located only to relative precision about
\(\sqrt{\varepsilon}\).
\end{theorem}

The reason is Taylor's theorem. Near the minimum,
\(f(x^* + h) \approx f(x^*) + \tfrac12 f''(x^*)\,h^2\). A displacement \(h\)
changes \(f\) by an amount proportional to \(h^2\). Once \(h^2\) falls below the
rounding level \(\varepsilon\,|f(x^*)|\), the function values at \(x^*\) and at
\(x^* + h\) are indistinguishable, and no method can tell which is lower. With
40 digits, \(\sqrt{\varepsilon}\) is about \(10^{-20}\), which is exactly the
agreement in the session. A root has no such limit, because \(\psi(x^* + h)\)
changes linearly in \(h\). This is why \B{FindRoot} on the derivative is the
better tool when a minimizer must be known to full precision. It is also why
the default \B{PrecisionGoal} of \B{FindMinimum} is half the working precision.

\calloutnote{Pitfall}{The value \(f(x^*)\) is accurate to nearly the full working
precision even when \(x^*\) is not, for the same reason. An error \(h\) in the
location costs only \(h^2\) in the value. A reported minimum value is usually
far more trustworthy than the reported location.}

\section{Descent in several variables}
\label{sec:opt-descent}

In several variables there is no bracket, and a method needs a direction as well
as a step. The gradient \(\nabla f\) points uphill, so \(-\nabla f\) is the
direction of steepest descent. \emph{Steepest descent} moves along it. It is a
poor method on anything but a round bowl. In a long narrow valley the gradient
points almost across the valley, and the iterates zigzag from wall to wall.

Newton's method uses curvature as well. Approximate \(f\) near the current
point \(x_k\) by its second-order Taylor model
\[
  m_k(p) = f(x_k) + \nabla f(x_k)^{\mathsf T} p + \tfrac12\, p^{\mathsf T} H_k\, p ,
\]
where \(H_k\) is the Hessian matrix of second derivatives. If \(H_k\) is
positive definite, the model has a unique minimum where its gradient vanishes,
at the \emph{Newton step}
\[
  p_k = -H_k^{-1} \nabla f(x_k).
\]
Near a minimum with a positive definite Hessian, Newton's method converges
quadratically. The number of correct digits roughly doubles at every step. Far
from the minimum \(H_k\) may be indefinite, and the Newton step may point
uphill.

\subsection{Line search}

Both steepest descent and Newton give a direction \(d_k\). A \emph{line
search}\index{line search} then chooses how far to go. Mathilda uses
backtracking with the \emph{Armijo condition}\index{Armijo condition}. Try a
step \(\alpha\), and accept it if
\[
  f(x_k + \alpha d_k) \le f(x_k) + c_1\, \alpha\, \nabla f(x_k)^{\mathsf T} d_k ,
  \qquad c_1 = 10^{-4}.
\]
Otherwise halve \(\alpha\) and try again. Because \(d_k\) is a descent
direction, the right-hand side decreases with \(\alpha\) at a rate proportional
to the slope, and a small enough step always satisfies the condition. The
condition asks for a decrease that is at least a small fraction of what the
linear model predicts. It rules out steps that reduce \(f\) by a negligible
amount.

\subsection{Quasi-Newton methods and BFGS}

Newton's method needs the Hessian, which costs \(n^2\) second derivatives and a
linear solve per step. A \emph{quasi-Newton method}\index{quasi-Newton method}
builds an approximation to the inverse Hessian from the gradients it has
already computed. After a step \(s_k = x_{k+1} - x_k\), the gradient changes by
\(y_k = \nabla f(x_{k+1}) - \nabla f(x_k)\). For a quadratic function,
\(y_k = H s_k\) exactly. So the new approximation \(B_{k+1}\) to the inverse
Hessian is required to satisfy the \emph{secant equation}
\[
  B_{k+1}\, y_k = s_k .
\]
Many matrices satisfy it. The \emph{BFGS update}\index{BFGS method} (Broyden,
Fletcher, Goldfarb, and Shanno, 1970) chooses the one closest to \(B_k\) in a
weighted norm, and keeps it symmetric and positive definite. With
\(\rho_k = 1/(y_k^{\mathsf T} s_k)\) it reads
\[
  B_{k+1} = B_k
    + \rho_k^2\,(s_k^{\mathsf T} y_k + y_k^{\mathsf T} B_k y_k)\, s_k s_k^{\mathsf T}
    - \rho_k\,\bigl(B_k y_k\, s_k^{\mathsf T} + s_k\, y_k^{\mathsf T} B_k\bigr).
\]
The search direction is \(d_k = -B_k \nabla f(x_k)\). The iteration starts with
\(B_0 = I\), so the first step is steepest descent, and the matrix learns the
curvature as the iterates move. BFGS converges superlinearly and needs only
gradients. It is the default method of \B{FindMinimum} in more than one
variable.

\calloutnote{Under the hood}{\mcode{fm\_run\_bfgs} in
\texttt{src/numerical\_calculus/findmin\_bfgs.c} applies exactly the update
above to a dense \(n \times n\) matrix. Two safeguards keep it well defined.
The update is skipped when \(y_k^{\mathsf T} s_k \le 10^{-12}\), because the
curvature condition \(y^{\mathsf T} s > 0\) is what keeps \(B_k\) positive
definite. If the direction \(-B_k\nabla f\) is ever not a descent direction,
\(B_k\) is reset to the identity. The iteration stops when
\(\|\nabla f\| < 10^{-\text{AccuracyGoal}}\) or when the step is smaller than
\(10^{-\text{PrecisionGoal}}\) relative to \(|x|\).}

Rosenbrock's function
\[
  f(x, y) = (1 - x)^2 + 100\,(y - x^2)^2
\]
is the standard test for these methods. Its minimum at \((1,1)\) lies at the
bottom of a long, curved, nearly flat valley. The usual starting point is
\((-1.2, 1)\), on the far side of the valley:

\begin{codepairs}
\pair{optimization/bfgs/1}{Name the objective once for the session.}
\pair{optimization/bfgs/2}{BFGS reaches \((1,1)\). The value
\(7\times10^{-21}\) is zero to machine precision.}
\pair{optimization/bfgs/3}{The gradient that BFGS uses is computed
symbolically by \B{D}, once, before the search starts.}
\pair{optimization/bfgs/4}{\mcode{StepMonitor} runs after every accepted step.
BFGS took 39 steps.}
\pair{optimization/bfgs/5}{\mcode{EvaluationMonitor} runs at every evaluation of
the objective or its gradient, including rejected line-search trials.}
\pair{optimization/bfgs/6}{\mcode{MaxIterations} stops the search early. After ten
steps the iterate is still on the far side of the valley.}
\pair{optimization/bfgs/7}{\mcode{Gradient} supplies the gradient explicitly. The
answer is the same.}
\pair{optimization/bfgs/8}{With bare variables, every starting value is \(1\).
Starting at the origin would put this objective on a saddle point, where the
gradient vanishes.}
\end{codepairs}

A monitor can also collect the iterates. Wrapping the call in \B{Reap} and
sowing \mcode{\{x, y\}} at each step records the whole path, which can then be
drawn over the level curves of the function:

\plotcell{optimization/rosenbrock-path}

The path shows how BFGS behaves. The first steps cross the valley floor
roughly, because \(B_0 = I\) knows nothing about the curvature. After a few
steps the matrix has learned the shape of the valley, and the iterates follow
the curved floor round to \((1, 1)\).

\calloutnote{Performance}{At machine precision \B{FindMinimum} compiles the
objective and each component of its symbolic gradient to bytecode
(\S\ref{sec:compile-auto}) and evaluates them on the register machine at every
trial point. The reference documentation records typical gains of about
\(3\times\) on a one-dimensional transcendental objective and \(4\times\) on
two-dimensional Rosenbrock. Setting \mcode{EvaluationMonitor} turns this off,
because the monitor has to run in the interpreter.}

\section{The local methods}
\label{sec:opt-local-methods}

BFGS is a good default, but no single method is best for every problem. The
\mcode{Method} option of \B{FindMinimum} selects from sixteen local methods.
Table~\ref{tab:opt-local} groups them by the information they use.

\begin{table}[ht]
\centering
\small
\begin{tabular}{@{}llll@{}}
\toprule
\mcode{Method} & uses & constraints & idea \\
\midrule
\mcode{"Brent"} & values & interval & golden section + parabolas (1 variable) \\
\mcode{"QuasiNewton"} & gradient & all smooth & dense BFGS \\
\mcode{"ConjugateGradient"} & gradient & all smooth & Polak--Ribi\`ere+ with restarts \\
\mcode{"LBFGSB"} & gradient & all smooth & BFGS from the last 10 steps \\
\mcode{"TNC"} & gradient & all smooth & truncated Newton, Hessian-free \\
\mcode{"Newton"} & Hessian & all smooth & Newton with modified Cholesky \\
\mcode{"NewtonCG"} & gradient & none & Newton step by conjugate gradients \\
\mcode{"Dogleg"} & Hessian & none & trust region, dogleg path \\
\mcode{"TrustNCG"} & gradient & none & trust region, Steihaug CG \\
\mcode{"TrustExact"} & Hessian & none & trust region, Mor\'e--Sorensen \\
\mcode{"TrustKrylov"} & gradient & none & trust region, Lanczos (GLTR) \\
\mcode{"Powell"} & values & boxes & conjugate directions \\
\mcode{"NelderMead"} & values & boxes & downhill simplex \\
\mcode{"SLSQP"} & gradient & all smooth & sequential quadratic programming \\
\mcode{"COBYLA"} & values & all & linear models in a trust region \\
\mcode{"COBYQA"} & values & all & quadratic models in a trust region \\
\bottomrule
\end{tabular}
\caption{The local methods of \mcode{FindMinimum}. ``All smooth'' means box
constraints plus smooth inequalities and equalities.}
\label{tab:opt-local}
\end{table}

\subsection{Newton, conjugate gradients, and limited memory}

\mcode{"Newton"} forms the exact Hessian by differentiating the symbolic
gradient. To survive an indefinite Hessian it factors \(H + \tau I\) by
Cholesky, increasing \(\tau\) until the factorization succeeds. The shifted
matrix is positive definite, so the step is always a descent direction, and
near a minimum \(\tau = 0\) works and full quadratic convergence returns.

\mcode{"ConjugateGradient"} needs no matrix at all. Each direction is the
negative gradient plus a multiple of the previous direction,
\(d_{k+1} = -g_{k+1} + \beta_k d_k\), with the Polak--Ribi\`ere+ choice
\(\beta_k = \max\bigl(0,\, g_{k+1}^{\mathsf T}(g_{k+1} - g_k)/\|g_k\|^2\bigr)\).
On a quadratic with exact line searches the directions are mutually conjugate
with respect to the Hessian, and the method terminates in at most \(n\) steps.

\mcode{"LBFGSB"} is the limited-memory form of BFGS. It never stores the
\(n\times n\) matrix. It keeps the ten most recent pairs \((s_i, y_i)\) and
computes \(B_k \nabla f\) directly by the two-loop recursion of Nocedal, in
\(O(10\,n)\) operations. The memory and the cost per step grow linearly with
\(n\), which is what makes quasi-Newton methods practical for thousands of
variables. \mcode{"TNC"} and \mcode{"NewtonCG"} take the same idea further.
They solve the Newton system \(H p = -g\) approximately by conjugate
gradients, and they touch \(H\) only through products \(Hv\), each estimated
from two gradients by the finite difference
\(Hv \approx (\nabla f(x + hv) - \nabla f(x))/h\).

\subsection{Trust regions}

A line search picks a direction first and a length second. A \emph{trust-region
method}\index{trust-region method} does the opposite. It fixes a radius
\(\Delta\) within which the quadratic model \(m_k\) is trusted, and chooses the
step by minimizing the model inside the ball:
\[
  \min_{\|p\| \le \Delta}\; \nabla f(x_k)^{\mathsf T} p + \tfrac12\, p^{\mathsf T} H_k\, p .
\]
The step is then judged by the ratio
\[
  \rho = \frac{f(x_k) - f(x_k + p)}{m_k(0) - m_k(p)}
\]
of the actual to the predicted decrease. If \(\rho > 0.15\) the step is
accepted. If \(\rho < 1/4\) the model was poor and \(\Delta\) shrinks by a
factor of four. If \(\rho > 3/4\) and the step reached the boundary,
\(\Delta\) doubles. These are SciPy's constants, and Mathilda uses them
unchanged. The four trust-region methods differ only in how they solve the
subproblem. \mcode{"Dogleg"} follows a two-segment path from the Cauchy point
to the Newton point. \mcode{"TrustNCG"} runs conjugate gradients and stops at
the boundary. \mcode{"TrustExact"} solves the subproblem almost exactly by the
Mor\'e--Sorensen iteration, which handles an indefinite Hessian.
\mcode{"TrustKrylov"} projects the problem onto a Krylov subspace by the
Lanczos process and solves the small tridiagonal problem exactly.

\subsection{Methods without derivatives}

Some objectives have no usable derivative. They may be noisy, piecewise, or
the output of a simulation. \mcode{"Powell"} minimizes along each of \(n\)
directions in turn with Brent's method, then replaces one direction by the net
displacement of the cycle. On a quadratic this builds up a set of conjugate
directions using function values alone.

\mcode{"NelderMead"} keeps a \emph{simplex}\index{Nelder--Mead simplex} of
\(n + 1\) points. Each step replaces the worst vertex \(x_h\). It first
reflects \(x_h\) through the centroid \(\bar x\) of the others, to
\(x_r = \bar x + (\bar x - x_h)\). If the reflected point is the best so far,
the simplex tries to expand further, to \(\bar x + 2(\bar x - x_h)\). If it is
still the worst, the simplex contracts halfway toward \(\bar x\), and if even
that fails, every vertex shrinks halfway toward the best one. The simplex
tumbles downhill, stretches along valleys, and contracts around a minimum. The
coefficients \(1, 2, \tfrac12, \tfrac12\) are the standard ones.

\begin{codepairs}
\pair{optimization/methods/1}{Rosenbrock's function again.}
\pair{optimization/methods/2}{Newton with the exact Hessian. It needed 20 steps
where BFGS needed 39.}
\pair{optimization/methods/3}{Limited-memory BFGS.}
\pair{optimization/methods/4}{A trust region with a near-exact subproblem
solve.}
\pair{optimization/methods/5}{The Nelder--Mead simplex uses no derivatives
at all, and it still reaches the minimum.}
\end{codepairs}

The iterates of four of these methods can be recorded with \mcode{StepMonitor},
exactly as before. The next figure plots \(\log_{10} f\) along each path. The
curves are, in order of their length, Newton (20 steps, orange), TrustExact
(29, red), LBFGSB (35, green), and QuasiNewton (39, blue).

\plotcell{optimization/convergence}

All four curves have the same shape. There is a long flat stretch while the
iterate travels along the valley, then a steep plunge once the method is close
enough for its model to be accurate. The plunge is the superlinear or
quadratic convergence of the theory. Most of the cost of a local method is in
the flat stretch, which is why a good starting point matters more than the
choice among good methods.

\subsection{Four minima, four starting points}

Himmelblau's function
\[
  h(x, y) = (x^2 + y - 11)^2 + (x + y^2 - 7)^2
\]
has four minima, each with value \(0\), and one local maximum. It shows the
local character of \B{FindMinimum} in two dimensions:

\begin{codepairs}
\pair{optimization/methods/6}{Name Himmelblau's function.}
\pair{optimization/methods/7}{From \((1,1)\) the answer is the integer minimum
\((3,2)\).}
\pair{optimization/methods/8}{From \((-1,1)\) the minimum in the second
quadrant.}
\pair{optimization/methods/9}{From \((-1,-1)\) the third.}
\pair{optimization/methods/10}{From \((1,-1)\) the fourth.}
\pair{optimization/methods/11}{\B{FindMaximum} from the origin finds the local
maximum, \(h \approx 181.617\).}
\end{codepairs}

\plotcell{optimization/himmelblau}

\subsection{A non-smooth objective}

The function \(g = |x - 1| + |y + 2| + |x - y|\) has corners along three lines.
Its minimum value is \(3\), attained on the whole segment \(x = y\),
\(-2 \le x \le 1\). The gradient methods still work here because the objective
is smooth almost everywhere, and the finite-difference fallback steps across
the kinks. Nelder--Mead makes no smoothness assumption at all:

\begin{codepairs}
\pair{optimization/nonsmooth/1}{A sum of absolute values.}
\pair{optimization/nonsmooth/2}{The default method lands on the segment of
minima, at \(x = y \approx 0.88\).}
\pair{optimization/nonsmooth/3}{Nelder--Mead lands at a different point with
the same value, because the minimizer is not unique.}
\end{codepairs}

\calloutnote{Pitfall}{A non-unique minimizer is common in practice, and
different methods (or different starting points) return different points with
the same minimum value. When you compare two answers to such a problem, compare
the values.}

\section{Constraints in a local search}
\label{sec:opt-local-constraints}

A constrained problem has the form
\[
  \min_x f(x) \quad \text{subject to} \quad g_i(x) \le 0,\; h_j(x) = 0 .
\]
\B{FindMinimum} accepts the constraints as the second element of
\mcode{\{f, cons\}}, written as ordinary comparisons joined by \mcode{\&\&}.
It sorts them into two kinds.

A comparison between one variable and a constant, such as \mcode{0 <= x <= 1},
is a \emph{box constraint}\index{box constraint}. Box constraints are enforced
exactly by \emph{projection}. After every trial step, each coordinate is clipped
back into its interval. The iterate never leaves the box.

Every other inequality or equality is a \emph{general constraint}, and the
default treatment is an \emph{augmented Lagrangian}\index{augmented Lagrangian}
method. The idea starts from the simpler \emph{quadratic
penalty}\index{penalty method}. Replace the constrained problem by the
unconstrained one
\[
  \min_x\; f(x) + \mu \Bigl( \sum_i \max(0, g_i(x))^2 + \sum_j h_j(x)^2 \Bigr),
\]
which charges a price for every violation. As \(\mu \to \infty\), the minimizer
of the penalized function approaches the constrained minimizer. The figure shows
this for the problem of minimizing \(x^2\) subject to \(x \ge 1\). The
penalized minimizer is \(\mu/(1 + \mu)\), which is \(0.5\), \(0.91\), and
\(0.99\) for \(\mu = 1, 10, 100\) (the blue, orange, and green curves).

\plotcell{optimization/penalty}

The figure also shows the weakness of the pure penalty. The minimizer
approaches the constraint only as \(\mu \to \infty\), and for large \(\mu\) the
penalized function has a steep wall that makes the unconstrained problem
badly conditioned. The augmented Lagrangian fixes this with Lagrange
multipliers. For each constraint \(c_k\) it keeps an estimate \(\lambda_k\) of
the multiplier and minimizes
\[
  f(x) + \sum_k \bigl(\lambda_k\, c_k(x) + \mu\, c_k(x)^2\bigr),
\]
where inactive inequalities are shifted out in the standard way. After each
inner minimization the multipliers are updated by the PHR rule (Powell,
Hestenes, and Rockafellar),
\[
  \lambda_k \leftarrow \lambda_k + 2\mu\, c_k(x)
  \qquad (\text{and } \lambda_k \leftarrow \max(0, \cdot) \text{ for an inequality}).
\]
The multiplier term shifts the minimizer onto the constraint at a finite
\(\mu\), so the method does not need the ill-conditioned limit. At the
solution, \(\lambda_k\) converges to the true Lagrange multiplier of the
Karush--Kuhn--Tucker conditions.

\calloutnote{Under the hood}{\mcode{fm\_run\_penalty} in
\texttt{src/numerical\_calculus/findmin\_penalty.c} starts with \(\mu = 1\)
and all \(\lambda_k = 0\), so its first round is exactly the quadratic penalty.
Each round runs the chosen inner solver (BFGS by default), updates the
multipliers, and multiplies \(\mu\) by ten, for at most nine rounds. It stops
when the iterate is feasible and has stopped moving between rounds. It also
remembers the best feasible iterate of the whole schedule. If a late,
ill-conditioned round drifts out of the feasible set, that earlier point is
returned instead.}

\begin{codepairs}
\pair{optimization/local-constraints/1}{Box constraints only. The unconstrained
minimum \((3,3)\) is outside the box, so the answer is the nearest corner.}
\pair{optimization/local-constraints/2}{The closest point of the unit disk to
\((2,1)\) is \((2,1)/\sqrt5\).}
\pair{optimization/local-constraints/3}{A linear program. The minimum of
\(x+y\) on this region is the corner \((7/3, 0)\).}
\pair{optimization/local-constraints/4}{SLSQP reaches the same corner with
\(y\) exactly \(0\).}
\pair{optimization/local-constraints/5}{An equality constraint: the minimum of
\(xy\) on the unit circle is \(-1/2\).}
\pair{optimization/local-constraints/6}{\B{FindMaximum} with a linear
constraint.}
\pair{optimization/local-constraints/7}{Problem~71 of Hock and Schittkowski,
with one inequality, one equality, and bounds. Its published optimum is
\(17.014\).}
\pair{optimization/local-constraints/8}{COBYQA solves the disk problem using
function values only.}
\end{codepairs}

\plotcell{optimization/disk}

The disk problem shows the geometry of a constrained minimum. The level curves
of the objective are circles around \((2,1)\). The solution is the point where
the smallest level curve that meets the feasible set touches it, and there the
gradient of the objective is parallel to the gradient of the constraint. That
parallel condition,
\(\nabla f(x^*) + \lambda\, \nabla g(x^*) = 0\) with \(\lambda \ge 0\), is the
Karush--Kuhn--Tucker condition for a single active inequality.

\subsection{Sequential quadratic programming}

The penalty approach treats constraints indirectly. \mcode{"SLSQP"} treats
them directly. At each iterate it replaces \(f\) by a quadratic model, whose
Hessian is a BFGS approximation to the Hessian of the Lagrangian, and
replaces each constraint by its linearization. The resulting \emph{quadratic
program}\index{quadratic programming} is solved exactly by the dual
active-set method of Goldfarb and Idnani, and the step is accepted by a line
search on an \(\ell_1\) penalty merit function. Near a solution SQP behaves
like Newton's method applied to the KKT conditions, and it converges
superlinearly. It is the method to reach for when a smooth constrained problem
has to be solved accurately, as with the HS71 problem above.

\mcode{"COBYLA"} and \mcode{"COBYQA"} are the derivative-free counterparts.
They build linear (COBYLA) or quadratic (COBYQA) models of the objective and
of every constraint by sampling points around the current iterate, and solve a
trust-region subproblem on the models. COBYLA's linear models are cheap and
robust to kinks, but they cannot follow a strongly curved valley. COBYQA's
quadratic models can.

\calloutnote{Pitfall}{Neither \mcode{"Powell"}, \mcode{"NelderMead"}, nor the
five trust-region methods accept general constraints. They reject such a
problem with the message \mcode{FindMinimum::nimpl}. Use the default method,
\mcode{"SLSQP"}, or one of the COBY methods instead.}

\section{Global search with NMinimize}
\label{sec:opt-global}

Rastrigin's function
\[
  r(x, y) = 20 + x^2 - 10\cos 2\pi x + y^2 - 10\cos 2\pi y
\]
is a paraboloid with a cosine ripple added. Every integer lattice point is near
a local minimum, and only the origin is global.

\plotcell[0.55\linewidth]{optimization/rastrigin}

A local method started almost anywhere on this surface stops in the nearest
dimple. \B{NMinimize} does not need a starting point:

\begin{codepairs}
\pair{optimization/nminimize/1}{Rastrigin's function in two variables.}
\pair{optimization/nminimize/2}{\B{FindMinimum} from \((2.3, -1.6)\) stops in
the dimple at \((1, -2)\), with value about \(5\).}
\pair{optimization/nminimize/3}{\B{NMinimize} finds the global minimum \(0\) at
the origin.}
\pair{optimization/nminimize/4}{Box constraints set the search region
explicitly.}
\pair{optimization/nminimize/5}{\B{SeedRandom} does not affect \B{NMinimize}.
This is the raw search result with no local polish.}
\pair{optimization/nminimize/6}{A different \B{SeedRandom} gives the
identical answer.}
\pair{optimization/nminimize/7}{The method option \mcode{"RandomSeed"}
changes the search, and the raw answer moves.}
\end{codepairs}

Two design decisions show in this session. First, the search is deterministic.
The stochastic engines draw from their own generator (SplitMix64) with a fixed
default seed, so the same call always gives the same answer, and a test that
checks it is stable. \B{SeedRandom}, which seeds \B{RandomReal} and the rest of
the system's random functions, has no effect on it. To vary the search, pass
\mcode{"RandomSeed" -> s} inside the \mcode{Method} option.

Second, the search needs a region. If a variable has a box constraint, the box
is the region. If it has none, \B{NMinimize} searches \([-10, 10]\) in that
coordinate. A variable with only one bound gets a span of \(20\) on the open
side. This default is why the first call above found the origin without being
told where to look.

\calloutnote{Pitfall}{A global optimizer with a bounded search region cannot
see a minimum outside it. If the true minimum of your problem might be far
from the origin, give the variables bounds. The one automatic exception,
described in \S\ref{sec:opt-global-constraints}, is a constrained problem
whose feasible set lies entirely outside the default region.}

\subsection{Differential evolution}

The default engine is \emph{differential evolution}\index{differential
evolution} (Storn and Price~\cite{stornprice}), in the variant called
DE/rand/1/bin. It maintains a population of \(N\) points
\(x_1, \dots, x_N\). In each generation every member \(x_i\) gets a challenger,
built in three steps.
\begin{enumerate}
\item \emph{Mutation.} Pick three other members \(x_{r_1}, x_{r_2}, x_{r_3}\) at
random and form the mutant
\[
  v = x_{r_1} + F\,(x_{r_2} - x_{r_3}), \qquad F = 0.6 .
\]
\item \emph{Crossover.} Build the trial point \(u\) coordinate by coordinate.
Take \(u_j = v_j\) with probability \(CR = 0.9\), and \(u_j = (x_i)_j\)
otherwise. One randomly chosen coordinate always comes from \(v\), so \(u\)
differs from \(x_i\).
\item \emph{Selection.} If \(u\) is better than \(x_i\), it replaces \(x_i\) in
the population.
\end{enumerate}

The mutation step is the clever part. The difference \(x_{r_2} - x_{r_3}\) of
two random members has the same distribution as the spread of the population.
Early on, when the population fills the search region, the steps are large and
the search is global. As the population gathers in the best basin, the
differences shrink and the search becomes local. The method adapts its own step
size with no tuning.

\calloutnote{Under the hood}{\mcode{nm\_de} in
\texttt{src/numerical\_calculus/nm\_de.c} uses a population of \(10n\) clamped
to \([15, 200]\). A mutant coordinate that leaves the box is not clipped to the
wall. It is reflected to a random point between its base vector and the
violated bound (``bounce-back''). Clipping piles members onto the boundary,
where their differences in that coordinate become zero and the coordinate
freezes. The run stops early once at least half the population is feasible and
the spread of their objective values is below
\(10^{-\text{AccuracyGoal}} + (1 + |f|)\,10^{-\text{PrecisionGoal}}\).}

After the last generation, \B{NMinimize} polishes. It takes the best
\(\min(2n, 50)\) distinct members of the final population, runs the local
solver of \S\ref{sec:opt-local-constraints} from each, and keeps the deepest
minimum. Polishing several members instead of only the best one matters on a
rugged surface. The best member of a half-converged population can sit in a
shallow basin next to a deeper one that another member has found. The polish
also explains the difference between the raw answers above, of order
\(10^{-10}\), and the polished answer \(0\).

The next figure shows the raw DE result on Rastrigin, with the polish switched
off, as a function of the generation budget \mcode{MaxIterations}. Each point is a
separate run with the same seed, so a larger budget continues the same
trajectory further. The vertical axis is \(\log_{10}\) of the best value.

\plotcell{optimization/de-budget}

The curve has two phases. For the first fifty or so generations it falls in
steps. Each flat stretch is the population exploring one basin, and each drop
is the moment a mutant lands in a better one. Just after generation fifty the
best value falls below \(1\), the height of the shallowest non-global minima,
so the population has reached the global basin. From then on the members
contract into it, and the best value falls by about a factor of ten every five
generations, reaching \(5\times10^{-10}\) at generation 100.

\calloutnote{Pitfall}{The generation budget depends on how the method is
named. With no \mcode{Method} option, DE runs up to \(150n\) generations. With
\mcode{Method -> "DifferentialEvolution"} written out, it runs up to \(100\).
Experiment~89 in \texttt{benchmarks/} found that this difference turns a
five-dimensional Rastrigin problem that the default solves into one that fails
on five of six seeds. If you name the method explicitly, set
\mcode{MaxIterations} as well.}

\section{A tour of the global engines}
\label{sec:opt-engines}

\B{NMinimize} has eight engines, and \mcode{Automatic} selects one of them. They share the variable handling, the
constraint treatment, the compiled objective, and the final polish, and differ
only in how they search. Table~\ref{tab:opt-global} summarizes them.

\begin{table}[ht]
\centering
\small
\begin{tabular}{@{}lll@{}}
\toprule
\mcode{Method} & search & reference \\
\midrule
\mcode{Automatic} & differential evolution, budget \(150n\) & \\
\mcode{"DifferentialEvolution"} & population, DE/rand/1/bin & Storn and Price 1997 \\
\mcode{"SimulatedAnnealing"} & \(\min(\max(2n,12),50)\) Metropolis chains & Kirkpatrick et al.\ 1983 \\
\mcode{"NelderMead"} & \(\min(2n,20)\) restarted simplexes & Nelder and Mead 1965 \\
\mcode{"RandomSearch"} & random starts, each polished & \\
\mcode{"BasinHopping"} & perturb, minimize, accept & Wales and Doye 1997 \\
\mcode{"DualAnnealing"} & generalized simulated annealing & Tsallis and Stariolo 1996 \\
\mcode{"SHGO"} & sample, graph, minimizer pool & Endres et al.\ 2018 \\
\mcode{"DIRECT"} & deterministic box subdivision & Jones et al.\ 1993 \\
\bottomrule
\end{tabular}
\caption{The global engines of \mcode{NMinimize}. Each accepts sub-options in
the form \mcode{Method -> \{"Name", "SearchPoints" -> k, \ldots\}}.}
\label{tab:opt-global}
\end{table}

\paragraph{Simulated annealing.} A single point walks at random. A move that
lowers the objective is always accepted. A move that raises it by \(\Delta\) is
accepted with the Metropolis probability \(e^{-\Delta/T}\), where the
temperature \(T\) decreases geometrically~\cite{kirkpatrick}. At high
temperature the walk crosses ridges freely, and as it cools it settles into a
basin. Mathilda runs \(\min(\max(2n, 12), 50)\) independent chains, polishes the
best point of each chain, and ranks the chains by their polished minima. It also
measures the typical size of \(\Delta\) near each chain's start, so that the
temperature is on the same scale as the objective.

\paragraph{Random search and Nelder--Mead.} \mcode{"RandomSearch"} draws random
points in the region and runs the local solver from each one. It is the
simplest global method, and it works only when a random start has a fair chance
of landing in the right basin. \mcode{"NelderMead"} runs the simplex method of
\S\ref{sec:opt-local-methods} from \(\min(2n, 20)\) random simplexes and keeps
the best polished result.

\paragraph{Basin hopping.} Each hop displaces the current point by a random
vector, then runs a full local minimization (the ``quench''). The Metropolis
test compares the two \emph{minimized} energies.
The walk therefore moves on the coarse landscape of basin floors, which is much
smoother than the original surface. Every fifty hops the step size is adjusted
toward an acceptance rate of one half.

\paragraph{Dual annealing.} Generalized simulated annealing draws its jumps from
a heavy-tailed Tsallis distribution with shape \(q_v = 2.62\), so that most
jumps are short but some cross the whole box. It accepts uphill moves by a
generalized Metropolis rule with shape \(q_a = -5\), starts at temperature
\(5230\), and reheats (``reanneals'') when the temperature falls below
\(2\times10^{-5}\) of the start. After each chain a local search refines the
best point, which is the second annealing that gives the method its name.
These are SciPy's defaults.

\paragraph{SHGO.} Simplicial homology global optimization samples the box,
joins the samples into a graph, and keeps the \emph{minimizer pool}: the samples
that are better than all their neighbours in the graph. Each pool vertex sits in
a different basin, so one local search from each pool vertex reaches every
basin the sampling resolved, with far fewer local searches than blind
multi-start.

\paragraph{DIRECT.} DIviding RECTangles (Jones, Perttunen, and
Stuckman~\cite{direct}) is deterministic. It scales the box to the unit cube and
samples its centre. At each iteration it selects the \emph{potentially optimal}
cells, those that could contain the minimum for some value of an unknown
Lipschitz constant, and splits each into thirds along its longest sides.
Large cells are split because the minimum might hide in them, and cells with
low centre values are split because the minimum is probably near them. The
default is the locally biased variant DIRECT-L of Gablonsky and Kelley.

All eight engines solve Rastrigin in two variables from their defaults, except
one. Differential evolution, Nelder--Mead, and DIRECT appeared above; here are
the rest:

\begin{codepairs}
\pair{optimization/engines/1}{Rastrigin's function again.}
\pair{optimization/engines/2}{Simulated annealing.}
\pair{optimization/engines/3}{Dual annealing.}
\pair{optimization/engines/4}{Basin hopping.}
\pair{optimization/engines/5}{SHGO lands exactly on the origin.}
\pair{optimization/engines/6}{Random search stops at a local minimum with value
about \(9.95\).}
\pair{optimization/engines/7}{With 400 starts instead of the default 16, it
finds the origin.}
\pair{optimization/engines/8}{Five seeds of random search give five answers,
none of them the global minimum.}
\pair{optimization/engines/9}{Five seeds of simulated annealing all find \(0\).}
\end{codepairs}

The random-search failure is instructive. Inside \([-10,10]^2\) Rastrigin has
about \(21^2 = 441\) basins, one near each integer point, and the global basin
is one of them. A random start lands in it with probability about \(1/441\),
so the default \(16\) starts succeed with probability about
\(1 - (440/441)^{16} \approx 4\%\). Four hundred starts succeed with
probability about \(60\%\). Multi-start search has no way to move
from one basin to the next. That is exactly what the population in DE, the hot
phase of annealing, and the hops of basin hopping supply.

\calloutnote{Theory}{The same arithmetic explains why random search fails in
high dimension. The reference documentation notes that the central basin of the
Griewank function on \([-600, 600]^{10}\) occupies about \(10^{-11}\) of the
box. No feasible number of random starts reaches it.}

\subsection{DIRECT and the lucky centre}

A deterministic method can look better than it is. DIRECT samples the centre of
the box first, and Rastrigin's minimum is at the centre of the symmetric box
\([-5.12, 5.12]^2\):

\begin{codepairs}
\pair{optimization/direct/1}{Rastrigin's function.}
\pair{optimization/direct/2}{Raw DIRECT on the symmetric box returns the origin
exactly, because the origin is its first sample.}
\pair{optimization/direct/3}{On an asymmetric box the raw DIRECT answer is close
to the origin, but not on it.}
\pair{optimization/direct/4}{The default polish finishes the job.}
\end{codepairs}

\calloutnote{Pitfall}{When you test a global optimizer, do not put the known
minimum at the centre of the search box. Test functions such as Rastrigin,
Ackley, and the sphere all have their minimum at the origin, and DIRECT (and
anything else that samples the centre) will find it for free.}

\subsection{Harder landscapes}

Rastrigin is regular and its global basin is central. The standard corpus of
test functions includes much harder cases. The next session tries four with
published optima: the six-hump camel (\(-1.0316\)), drop-wave (\(-1\) at the
origin), Schwefel's function (\(0\) at \(x_i = 420.9687\), near the corner of
the box), and the Eggholder (\(-959.6407\) at \((512, 404.2319)\), on the edge
of the box).

\begin{codepairs}
\pair{optimization/hard/1}{The six-hump camel function. The answer is one of its
two symmetric global minima.}
\pair{optimization/hard/2}{Drop-wave: global value \(-1\) at the origin.}
\pair{optimization/hard/3}{Schwefel's function in two variables. The residual
\(2.5\times10^{-5}\) comes from the rounded constant \(418.9829\).}
\pair{optimization/hard/4}{The Eggholder function.}
\pair{optimization/hard/5}{Default DE stops in a deep local basin at
\(-935.34\).}
\pair{optimization/hard/6}{Simulated annealing reaches the global basin, within
\(0.004\) of the published optimum.}
\end{codepairs}

\plotcell[0.6\linewidth]{optimization/eggholder}

The Eggholder shows that a global optimizer's answer is a claim. It is never a
proof.
Its surface is a set of diagonal ridges and troughs, and its best trough runs
into the edge of the box. Differential evolution converged on the white point,
a trough nearly as deep as the global one and far from it. Simulated annealing,
with its many independent chains, found the red point on the edge. The
reference documentation records that Mathematica's default \mcode{NMinimize}
stops at the same \(-935.34\).

\section{Constraints in a global search}
\label{sec:opt-global-constraints}

Inside the global search, Mathilda handles general constraints with
\emph{Deb's feasibility rules}\index{Deb's feasibility rules} instead of a
weighted penalty. To compare two candidate points:
\begin{enumerate}
\item a feasible point beats an infeasible one;
\item of two feasible points, the one with the smaller objective wins;
\item of two infeasible points, the one with the smaller total violation wins.
\end{enumerate}
The rules need no penalty weight, so there is nothing to tune. The objective
and the constraints are never added together, so a large objective cannot
drown a small violation or the reverse. Box constraints define the search
region and are enforced by projection, as before. After the search, the
polish uses the augmented Lagrangian of \S\ref{sec:opt-local-constraints}.

The total violation is \(\sum_i \max(0, g_i)^2 + \sum_j h_j^2\), and a point
counts as feasible when this is below a threshold. How large that threshold
should be turns out to be a delicate question, and getting it wrong produced
a real bug.

\subsection{Two thresholds for two questions}

In August 2026, while building the benchmarks of experiment~89, the author
found that
\[
  \mcode{NMinimize[\{x\^{}2 + y\^{}2, x + y >= 2\}, \{x, y\}]}
\]
returned \(1.9998\) at \(x = 0.999971\), \(y = 0.999929\). The sum \(x + y\) is
\(1.99990\), so the returned point violated its own constraint by
\(10^{-4}\). No warning was printed. SciPy and Mathematica both returned the
feasible point \((1,1)\) with value \(2\).

The cause was one constant doing two jobs. The feasibility test compared the
\emph{squared} violation total with \(10^{-8}\), which is a tolerance of
\(10^{-4}\) on the violation itself. The same constant served two purposes. It
ranked candidates during the search, and it decided whether the final answer
was feasible enough to return. Differential evolution stopped improving the
constraint as soon as the squared violation crossed \(10^{-8}\), and the
return test then accepted the result.

The two purposes need different thresholds. Deb's rules consult the objective
only when \emph{both} points count as feasible. If the ranking threshold is
tighter than the residual the search can actually reach, no point ever counts
as feasible. The rules then compare violations only, and the search stops
optimizing the objective altogether. This was measured. A tight ranking
threshold sent a 15-dimensional minimax test problem from its optimum
\(0.125116\) to \(1.85479\). The return test answers a different question. It
certifies to the caller that the point satisfies the constraints, and that
certificate should be tight.

The fix split the constant in two, each named for its question and stated as a
violation that is squared where it is defined:
\begin{itemize}
\item \mcode{NM\_FEAS\_RANK}, a violation of \(10^{-4}\), used only to rank
points during the search. It is loose on purpose, and bit-identical to the old
constant, so every search trajectory is unchanged.
\item \mcode{NM\_FEAS\_RETURN}, a violation of \(10^{-5}\), used for the
post-polish selection and the final decision. It is enforced. A result that
cannot meet it is reported as \mcode{\{Infinity, x -> Indeterminate\}}, the
same form \B{NMinimize} uses for an infeasible problem.
\end{itemize}
The return bound is \(10^{-5}\) and not tighter because it is set where every
supported class of problem can meet it. The measured residuals were about
\(4\times10^{-12}\) for continuous problems with seven of the eight methods
(Nelder--Mead reached \(10^{-6}\)), and about \(6\times10^{-6}\) for
mixed-integer problems, whose continuous coordinates are
refined with the integers pinned. A tighter bound would turn solvable problems
into \mcode{Infinity}.

\begin{codepairs}
\pair{optimization/global-constraints/1}{The problem that exposed the bug. The
answer is now exactly \((1,1)\).}
\pair{optimization/global-constraints/2}{The equality version.}
\pair{optimization/global-constraints/3}{The disk and the half-plane
\(x + y \ge 3\) do not meet. An empty feasible set gives \mcode{Infinity}.}
\pair{optimization/global-constraints/4}{A disjunction: \(x \le -2\) or
\(x \ge 2\). The minimum of \(x^2\) is \(4\), on the boundary of one branch.}
\pair{optimization/global-constraints/5}{Himmelblau's function restricted to two
small disks, one around each of two minima. The answer is the minimum inside
one of them.}
\pair{optimization/global-constraints/6}{The feasible set \(x + y \ge 80\) lies
entirely outside the default region \([-10,10]^2\). The region grows until it
finds it.}
\end{codepairs}

\calloutnote{Under the hood}{The two predicates are \mcode{nm\_better} and
\mcode{nm\_better\_return} in
\texttt{src/numerical\_calculus/findmin\_nm\_common.c}. Both are Deb's rules
with a different threshold. The long comment above the two constants in
\texttt{findmin\_internal.h} records the history, including the value
\(10^{-6}\) that was tried first for the return bound and rejected because it
turned a mixed-integer problem with the exact solution \((1,1)\) into
\mcode{Infinity}.}

A disjunction \mcode{c1 || c2} is feasible when at least one branch holds.
\B{NMinimize} scores it by the smallest of the branch violations, which is zero
exactly when some branch is satisfied. The minimum is not differentiable, but
the global engines use only function values, so they consume it directly.
\B{FindMinimum}, whose penalty method needs a smooth penalty, rejects \mcode{Or}.

The last example shows \emph{region expansion}\index{region expansion}. When the
default region contains no feasible point, the unbounded coordinates grow by
powers of ten, up to \(\pm10^5\), and the search is repeated. The first attempt
always uses the base region and the base seed, so a problem that is feasible
there is solved exactly as before.

\section{Integer and linear problems}
\label{sec:opt-integer}

Mathilda has no separate linear-programming function. There is no
\mcode{LinearOptimization}, \mcode{LinearProgramming}, or
\mcode{QuadraticOptimization}, and \mcode{FindMinimum} rejects the method
names \mcode{"InteriorPoint"} and \mcode{"LinearProgramming"} as not yet
implemented. A linear or quadratic program is posed to \B{NMinimize} or
\B{NMaximize} like any other constrained problem, and a variable is made
integer by the constraint \mcode{Element[x, Integers]}.

An integer program shows how much integrality changes a problem. Maximize \(x + y\)
subject to \(-2x + 2y \ge 1\), \(-8x + 10y \le 13\), and \(x, y \ge 0\):

\begin{codepairs}
\pair{optimization/integers/1}{The linear relaxation. Its optimum is the vertex
\((4, 4.5)\), with value \(8.5\).}
\pair{optimization/integers/2}{With integer variables the optimum is
\((1, 2)\), with value \(3\). Integer answers print as exact integers.}
\end{codepairs}

\plotcell[0.5\linewidth]{optimization/ilp}

The feasible region of the relaxation is the thin shaded triangle. Rounding its
optimum \((4, 4.5)\) gives \((4, 4)\) or \((4, 5)\), and neither is feasible.
The only integer points in the triangle are \((0,1)\) and \((1,2)\), far from the
relaxed optimum. In general the integer optimum can be arbitrarily far from the
relaxed one, which is why integer programming is NP-hard while linear
programming is solvable in polynomial time.

\calloutnote{Under the hood}{Integer coordinates are rounded whenever a
candidate is generated, so the global search moves on the lattice. The polish
for an integer problem is a coordinate descent that tries steps of \(\pm1\) and
\(\pm2\) on each integer coordinate and accepts any improvement under Deb's
rules. A swap move handles ``one-hot'' groups, equalities of the form
\(x_1 + \dots + x_m = k\) over 0--1 variables such as the rows of an assignment
matrix. When continuous variables are also present, the polish solves the
continuous relaxation, rounds, and refines the continuous coordinates with the
integers fixed.}

\begin{codepairs}
\pair{optimization/integers/3}{The 0--1 knapsack. Three items of weight 10, 20,
30 and value 60, 100, 120 go in a bag of capacity 50. The best choice is the
second and third items.}
\pair{optimization/integers/4}{The relaxation takes a fraction of the third
item and is worth 240.}
\pair{optimization/integers/5}{A mixed problem, with \(x\) integer and \(y\)
real, and an equality. The exact solution is \(x = y = 1\), value \(3\). The
continuous coordinate is refined to about \(6\times10^{-6}\), inside the
return bound.}
\pair{optimization/integers/6}{The optimum \((15, 3)\) lies outside the default
region. The relaxation step of the polish finds it.}
\end{codepairs}

\calloutnote{Tip}{Write the integer declaration inside the constraints, joined
with \mcode{\&\&}. When the problem is a list of three or more elements,
\mcode{\{f, c1, c2\}}, both \B{NMinimize} and \B{NMaximize} treat the extra
elements as further constraints, implicitly \mcode{And}-ed. (An earlier
\B{NMaximize} negated the whole list, which threaded \mcode{Times[-1, \dots]}
over the constraints and returned a meaningless value; it now negates only the
objective. The form \mcode{\{f, c1 \&\& c2\}} works too.)}

\section{Many variables}
\label{sec:opt-many}

Test problems are usually stated for \(n\) variables, and writing
\(x_1, \dots, x_{10}\) by hand is tedious. \B{NMinimize} accepts
\emph{indexed variables} such as \mcode{x[i]}, generated by \B{Table} or
\B{Array}, and so does \B{FindMinimum}: each accessor is rewritten to a fresh
scalar symbol and the result is reported over the original \mcode{x[i]}. You can
also generate plain symbols with \B{Symbol} and optimize over those:

\begin{codepairs}
\pair{optimization/indexed/1}{Rastrigin's function in five variables, with
indexed variables.}
\pair{optimization/indexed/2}{Rosenbrock's function in five variables.
\B{Array} builds the variable list.}
\pair{optimization/indexed/3}{Six generated symbols for \B{FindMinimum}.}
\pair{optimization/indexed/4}{The six-variable Rosenbrock function.}
\pair{optimization/indexed/5}{\B{FindMinimum} holds its arguments, so the
objective and the variable specification are wrapped in \B{Evaluate}.}
\end{codepairs}

\calloutnote{Performance}{Indexed variables once carried a large penalty:
experiment~89 measured the same five-dimensional Rastrigin problem at 2.9~ms with
explicit variables and 119.7~ms with \mcode{Table[v[i], \{i, 1, 5\}]}, a factor
of 41. The cause was never the optimizer --- its rewrite to fresh symbols and the
compiled fast path already handled indexed variables --- but \B{Sum}'s Gosper
stage churning \B{Simplify} on the opaque-indexed summand
(\mcode{v[i]\^{}2 - 10 Cos[2 Pi v[i]]}) before failing. With a structural guard in
that stage the two forms now run at parity, so indexed variables are a
convenience with no speed cost.}

\section{How fast, and how reliable}
\label{sec:opt-performance}

A global optimizer evaluates its objective thousands of times, so the cost of
one evaluation dominates everything else. Before the search begins,
\B{NMinimize} compiles the objective and each general constraint to bytecode
(Chapter~\ref{ch:compilation}) and evaluates them on the register machine at
every trial point. If an expression cannot be compiled, or the compiled code
reports a domain error or a non-finite result at some point, that point falls
back to the interpreter, so compilation changes the speed and never the answer.
SciPy's optimizers call back into a Python function for every evaluation.

Experiment~89 in \texttt{benchmarks/} races \B{NMinimize} and \B{NMaximize}
against SciPy~1.17.1 and Mathematica. Table~\ref{tab:opt-bench} gives a
selection. Times are the best of three runs after a warm-up.

\begin{table}[ht]
\centering
\small
\begin{tabular}{@{}lrrr@{}}
\toprule
Case & Mathilda & Mathematica & SciPy \\
\midrule
Rastrigin 2-D, DE & 0.114 & 213.8 & 24.9 \\
Ackley 10-D, DE & 10.04 & 1188 & 742.3 \\
Rosenbrock 5-D, DE & 1.529 & 916.6 & 536.2 \\
Beale 2-D, default & 0.094 & 4.333 & 37.7 \\
inequality-constrained & 0.255 & 0.941 & 0.189 \\
equality-constrained & 0.218 & 0.414 & 0.186 \\
mixed integer & 0.019 & 0.646 & 5.540 \\
Rastrigin 5-D, explicit variables & 2.933 & 740.6 & 163.3 \\
Rastrigin 5-D, indexed variables & 2.98 & 773.4 & 162.2 \\
\bottomrule
\end{tabular}
\caption{Selected timings from experiment 89, in milliseconds. The indexed-variable
row is shown after the \B{Sum}/Gosper fix, now at parity with explicit variables.}
\label{tab:opt-bench}
\end{table}

Across its 18 cases, experiment~89 finds Mathilda between \(1.4\times\) and
\(350\times\) faster than the faster of the two competitors. The two
constrained rows are the exception. SciPy's SLSQP is slightly faster there, by
factors of \(1.35\) and \(1.17\). Those two rows were also the ones that
exposed the feasibility bug of \S\ref{sec:opt-global-constraints}. After the
fix the inequality case takes \(0.259\)~ms against the \(0.255\)~ms in the
table, so correctness cost nothing measurable. The indexed-variable row once
showed a \(41\times\) penalty; with the \B{Sum}/Gosper fix of \S\ref{sec:opt-many}
it is at parity with the explicit row, and Mathilda keeps its \(56\times\) lead
there too.

The races of the individual engines tell the same story. Against SciPy's
implementations of the same algorithms with the same parameters, dual
annealing (experiment~81) is about \(100\)--\(600\times\) faster, basin hopping
(experiment~85) about \(100\)--\(650\times\), and DIRECT (experiment~83), whose
SciPy version is itself compiled C, \(1.5\)--\(20\times\). The reference
documentation reports the local constrained solvers against SciPy at about
\(17\times\) for SLSQP on HS71, \(40\)--\(200\times\) for COBYLA, and
\(30\)--\(390\times\) for COBYQA. Local minimization of a cheap function is a
different matter. On two-dimensional Rosenbrock, the weekly benchmark report
(\texttt{benchmarks/REPORT.md}) has \B{FindMinimum} at \(0.190\)~ms, SciPy at
\(3.6\)~ms, and Mathematica at \(0.158\)~ms, so Mathilda is \(1.2\times\) behind
Mathematica there.

\subsection{Speed is not robustness}

Experiment~90 is the companion of experiment~89, and it asks a different
question. Did the optimizer find the published global optimum at all? It uses
the landscapes that experiment~89 excluded because the systems disagreed. On
seven hard functions Mathilda's default now solves four, matching SciPy's
differential evolution, with a different set: Griewank in five dimensions,
Rastrigin in ten, Styblinski--Tang, and drop-wave, which SciPy misses.

What closed the gap was not a bigger budget but a stronger default search. Under
\mcode{Method -> Automatic} the differential-evolution engine now seeds its
population by Latin-hypercube sampling (even coverage of each coordinate, which
is what let Rastrigin's separable basin be assembled by crossover), mutates
toward the running best (current-to-best/1), dithers its scaling factor each
generation, carries SciPy's larger \(15d\) population, and keeps the best of a few
independent runs so a deceptive low-dimensional basin is not missed by seed luck.
Naming \mcode{"DifferentialEvolution"} explicitly keeps the classic configuration
unchanged.

The three remaining misses are honest limitations rather than failures of
persistence. Schwefel and the Eggholder are \emph{unbounded as written} --- their
objectives descend below the published optimum once the variables leave the
intended domain, so returning that optimum without bounds would be wrong, not
right. SciPy is handed the standard box bounds; given the same bounds Mathilda
solves both (the Eggholder to \(-959.641\)), which takes the count to six of
seven. Only Bukin~N.6, whose minimum lies in a razor-thin valley, defeats plain
differential evolution even with bounds, as it does SciPy. On a hard problem,
then, the first move is to state the variable bounds; after that, raise the budget
with \mcode{MaxIterations} or \mcode{"SearchPoints"}, try more than one engine,
and try more than one \mcode{"RandomSeed"}. If the answers agree, you have
evidence. If they do not, you have learned something important about your problem.

\section{Testing an optimizer}
\label{sec:opt-testing}

The feasibility bug of \S\ref{sec:opt-global-constraints} shipped even though
the test suite, \texttt{tests/test\_nminimize.c}, had 29 tests that all passed,
several of them constrained. Every one of those tests checked the objective to
within \(10^{-3}\), and \(1.9998\) is within \(10^{-3}\) of \(2\). None checked
that the returned point satisfied the constraint. After the fix the author
audited the whole suite, which by then had 83 tests, in two ways.

The structural audit classified every test by what it asserts. Thirty-one
constrained tests never checked their constraints at all. Seventeen
feasibility checks were one-sided, of the form ``\(x + y \le 2 + \epsilon\)''.
A one-sided check catches a point that overshoots the boundary and is blind to
one that falls short of it, and falling short is exactly what the bug did.

The mutation audit asked how many tests would notice a wrong answer. A
temporary change made the result builder return a deliberately wrong answer,
and a non-fatal assertion macro recorded the pass/fail map of the whole suite
in one run. With the returned point 10\% wrong and the objective value left
correct, 26 of 83 tests failed. With the objective 10\% wrong and the point
left correct, 60 of 83 failed. The suite was more than twice as good at
noticing a wrong number as a wrong answer. After the audit added two-sided
feasibility assertions to the eleven tests with real constraints, and
documented why each remaining constrained test has none, a 10\%-wrong point was
caught by 32 of 83.

The lesson generalizes to any numerical code, and the audit wrote it into the
header of the test file as a policy. An objective tolerance may be loose,
because a stochastic search does not reach the optimum exactly. A feasibility
tolerance must be tight and two-sided, because feasibility is what makes the
answer an answer. Where a feasibility tolerance has to be loose, that is a
finding about the solver, and it should be recorded as one.

\calloutnote{Pitfall}{The same discipline applies when you check an optimizer's
answer yourself. Substitute the returned point into every constraint, as pair~8
of the first session did for the objective. A minimum value that looks right
proves nothing about feasibility.}

\section{Featured examples}
\label{sec:opt-featured}

The examples in this section are short. Each poses a familiar problem, solves it
in a line or two, and draws the answer. They are chosen so that you can check
every result by hand or against a known formula.

\subsection*{Design the cheapest can}

A drinks can must hold 355~ml. Which radius and height use the least metal? The
surface area of a closed cylinder is \(2\pi r^2 + 2\pi r h\), and the volume
constraint is \(\pi r^2 h = 355\). Calculus gives \(h = 2r\) and
\(r = (355/2\pi)^{1/3} \approx 3.837\)~cm, so the ideal can is as tall as it is
wide. The plot shows the area with \(h\) eliminated, as a function of \(r\).

\plotcell[0.55\linewidth]{optimization/can}

Real cans are taller than this. One reason is that the top and bottom are made
of thicker metal than the wall. Weighting the two terms of the area differently is a
one-character change to the objective.

\subsection*{Fold the largest open box}

Cut a square of side \(x\) from each corner of a \(30 \times 20\) sheet and fold
up the sides. The volume is \((30 - 2x)(20 - 2x)x\). Setting its derivative to
zero gives a quadratic whose smaller root is
\(x = (50 - 10\sqrt7)/6 \approx 3.924\).

\plotcell[0.55\linewidth]{optimization/openbox}

The constraint \(0 \le x \le 10\) keeps the box physical. Without it the cubic
grows without bound for large \(x\).

\subsection*{Fit a curve to data}

Least-squares fitting is minimization. Twenty-five noisy measurements of an
exponential decay are fitted by the model \(a e^{-kt} + c\) by minimizing the
sum of squared residuals. The data were generated with \(a = 3\), \(k = 0.7\),
\(c = 0.5\) and noise of at most \(0.1\), and the fit recovers all three to
within a few percent.

\plotcell[0.6\linewidth]{optimization/fit}

\B{Fit} solves \emph{linear} least-squares problems directly. A model that is
nonlinear in its parameters, like this one, needs an iterative minimizer.

\subsection*{Place a warehouse to minimize distance}

Five stores sit at known positions, and their weekly demand is \(3, 1, 2, 1,
1\) truckloads. Where should the warehouse go to minimize total truck travel?
The objective is the demand-weighted sum of straight-line distances, the
Fermat--Weber problem. It has no closed-form solution, and the objective is not
differentiable at the stores.

\plotcell[0.55\linewidth]{optimization/warehouse}

The disk sizes show the demand. The warehouse (red) is pulled toward the
heaviest store at the origin. At the optimum the weighted unit vectors toward
the five stores sum to zero, a balance of forces that Varignon demonstrated
with strings and weights through holes in a table.

\subsection*{Schedule production under constraints}

A plant makes two products earning 3 and 5 units of profit per batch. Three
workshops limit the batches: \(x \le 4\), \(2y \le 12\), and \(3x + 2y \le
18\). This is the Wyndor Glass problem of Hillier and Lieberman's textbook,
with optimum \(36\) at \((2, 6)\).

\plotcell[0.5\linewidth]{optimization/production}

The shaded pentagon is the set of feasible plans, and the grey lines are two
level lines of the profit, \(3x + 5y = 24\) and \(3x + 5y = 36\). A linear
objective is maximized at a vertex of the polygon, which is the basis of the
simplex method.

\subsection*{Build an efficient portfolio}

Three assets have expected returns \(6\%\), \(10\%\), and \(14\%\), and a known
covariance matrix. For a required return \(r\), the least risky long-only
portfolio minimizes the variance \(w^{\mathsf T} C w\) subject to
\(\sum w_i = 1\), \(\mu^{\mathsf T} w = r\), and \(w \ge 0\). Solving this
quadratic program for a range of \(r\) traces Markowitz's \emph{efficient
frontier}\index{efficient frontier}.

\plotcell[0.55\linewidth]{optimization/portfolio}

The horizontal axis is the standard deviation of the portfolio and the vertical
axis its expected return. Each point is one SLSQP solve. The curve bends back
at the bottom. The least risky portfolio mixes all three assets and is less
risky than the safest single asset, whose standard deviation is \(0.1\).

\subsection*{Reach a target with a robot arm}

A planar arm has links of length \(1\) and \(0.7\). Which joint angles put the
hand at \((1.2, 0.8)\)? Inverse kinematics is a system of two equations in two
unknowns, and minimizing the squared distance from hand to target solves it.

\plotcell[0.5\linewidth]{optimization/robot}

The residual \(5\times10^{-19}\) says the hand is on the target. A second
solution, with the elbow bent the other way, is reached from a different
starting angle. For a redundant arm with more joints than equations, the same
formulation still works, and a second term in the objective can choose among
the solutions.

\subsection*{Follow a ray of light}

Fermat's principle says that light travels between two points along the path
of least time. A ray goes from \((0, 1)\) in air, at speed \(1\), to
\((2, -1)\) in glass, at speed \(2/3\), crossing the surface \(y = 0\) at some
point \((x, 0)\).

\plotcell[0.5\linewidth]{optimization/snell}

The ratio of the sines of the angles to the normal is exactly \(1.5\), the
ratio of the two speeds. Minimizing travel time has derived Snell's law. The
straight line (blue) is shorter, but the ray (red) spends less time in
the slow glass.

\subsection*{Hang a chain}

A chain of ten links, each of length \(1/4\), hangs between \((0, 0)\) and
\((2, 0)\). It settles where its potential energy, the sum of the heights of its
joints, is least, subject to one equality constraint per link. The problem has
eighteen variables and ten nonlinear equality constraints, and the default
augmented Lagrangian method solves it.

\plotcell[0.55\linewidth]{optimization/chain}

The shape is a discrete catenary. As the number of links grows it converges to
the curve \(y = a\cosh((x - 1)/a) + b\) of the calculus of variations.

\subsection*{Find the lowest point of a rugged landscape}

The Ackley function is nearly flat far from the origin, covered with small
dimples, and has one deep funnel at the centre. Its global minimum is \(0\) at
the origin. The default call finds it without any bounds.

\plotcell[0.6\linewidth]{optimization/ackley}

The dimples trap local methods and the flat outer region gives a gradient
almost no information. Differential evolution needs neither. Its population
samples the whole region, and its difference vectors shrink as the members
collect in the funnel.

\subsection*{Where this connects}

Optimization sits on top of most of the rest of the system. The symbolic
gradients and Hessians that drive BFGS, Newton, and SLSQP come from the
differentiation of \S\ref{sec:numcalc} and its exact neighbours, and the
compiled evaluation that makes every search fast is the machinery of
Chapter~\ref{ch:compilation}. \B{FindRoot} and \B{FindMinimum} are two views of
the same problem, as the Gamma example showed, and a minimizer located to full
precision is often best computed as a root of the derivative. The featured
examples drew on the graphics of Chapter~\ref{ch:graphics}. Two gaps remain
that a reader coming from Mathematica will notice. There is no dedicated
linear or quadratic programming function, so such problems go to
\B{NMinimize}, and there are no \mcode{FindArgMin} or \mcode{NArgMin} shortcuts,
so the location is read from the rules of the full result. For the complete
list of methods and sub-options, type \texttt{?FindMinimum} or
\texttt{?NMinimize} at the prompt.
